% =============================================================================
\section*{Synthèse exécutive}
\addcontentsline{toc}{section}{Synthèse exécutive}
% =============================================================================

Ce rapport présente la conception et l'évaluation d'une métaheuristique pour
l'optimisation de trajectoires 2D modélisées par des courbes de Bézier en
présence d'obstacles.  L'objectif est de minimiser une fonction de coût
composite (longueur, collisions, courbure, sortie de zone) sous contrainte de
temps (60~secondes), la qualité étant évaluée en moyenne sur 10~exécutions.

\paragraph{Phase exploratoire.}
Nous avons implémenté et évalué 19~algorithmes et variantes, couvrant quatre
familles : algorithmes génétiques (GA), évolution différentielle (DE, jDE,
SHADE, L-SHADE), stratégies d'évolution (CMA-ES, IPOP, BIPOP, LMCMA), et
optimisation par essaim particulaire (PSO, CLPSO, FPSO).  Des mécanismes
complémentaires ont été testés : hybridation avec A*, modèles de substitution
(\textit{surrogate}), gestion avancée des contraintes (Lagrangien augmenté,
\textit{stochastic ranking}), et bornes étendues.

\paragraph{Résultats de la phase exploratoire.}
L'analyse expérimentale sur quatre instances de difficulté croissante a mis en
évidence la supériorité de CMA-ES sur les instances corrélées et le rôle
décisif des bornes étendues pour l'instance labyrinthique \texttt{prob4}.

\paragraph{Algorithme final : OptiPath.}
Sur la base de ces observations, nous avons conçu \textbf{OptiPath}, une
métaheuristique basée sur CMA-ES combinant trois mécanismes :
\begin{enumerate}[nosep]
  \item \textbf{Seeding massif} ($\sim$500~solutions initiales) suivi d'une
        optimisation locale sur les 15~meilleures, pour identifier les bassins
        d'attraction prometteurs.
  \item \textbf{Portfolio à deux marges} : deux instances CMA-ES/BIPOP
        fonctionnent en alternance avec des bornes étendues différentes
        (marge~5 pour les instances simples, marge~15 pour les labyrinthes).
  \item \textbf{Restart adaptatif} : si la meilleure solution dépasse un seuil
        de coût après 5~secondes, les deux CMA-ES sont relancés depuis des
        seeds diversifiés, offrant plusieurs tentatives indépendantes pour
        trouver le bassin optimal.
\end{enumerate}

\paragraph{Résultats finaux (moyenne sur 10~exécutions).}
OptiPath obtient : \texttt{prob1}~$= 36{,}80$ ; \texttt{prob2}~$= 31{,}42$ ;
\texttt{prob3}~$= 29{,}02$ ; \texttt{prob4}~$= 94{,}60$.  Sur l'instance
labyrinthique \texttt{prob4}, OptiPath réduit le coût d'un facteur~16 par
rapport aux meilleures méthodes de la phase exploratoire ($94{,}60$ contre
$\sim$1\,517 pour L-SHADE avec bornes étendues).
